\subsection{Continuous linear representations}

Let \(\mathbf{G}\) be a topological group, \(\mathbf{E}\) a locally convex space, \(U\) a linear representation of \(\mathbf{G}\) in \(\mathbf{E}\) .

DEFINITION 1. --- (i) U is said to be separately continuous if, for every \(s \in G\) , \(U(s)\) is a continuous endomorphism of E, and if, for every \(x \in E\) , the mapping \(s \mapsto U(s)x\) of G into E is continuous.

\begin{enumerate}
\def\labelenumi{(\roman{enumi})}
\setcounter{enumi}{1}
\item
  \(U\) is said to be continuous if \((s,x)\mapsto U(s)x\) is a continuous mapping of \(\mathbf{G}\times \mathbf{E}\) into \(\mathbf{E}\) .
\item
  \(U\) is said to be equicontinuous if it is continuous and if the set of endomorphisms \(U(s)\) , where \(s\) runs over \(G\) , is equicontinuous.
\end{enumerate}

Remarks. --- 1) To say that U is separately continuous means that \(s \mapsto U(s)\) is a continuous mapping of G into the space \(\mathcal{L}(\mathrm{E};\mathrm{E})\) of continuous endomorphisms of E, equipped with the topology of pointwise convergence.

\begin{enumerate}
\def\labelenumi{\arabic{enumi})}
\setcounter{enumi}{1}
\tightlist
\item
  To say that \(U\) is continuous is equivalent to the following set of three conditions:
\end{enumerate}

\begin{enumerate}
\def\labelenumi{\alph{enumi})}
\tightlist
\item
  for every \(s \in G\) , \(U(s)\) is continuous; b) there exists a neighborhood V of e such that \(U(V)\) is equicontinuous; c) there exists a total set D in E such that, for every \(x \in D\) , the mapping \(s \mapsto U(s)x\) is continuous.
\end{enumerate}

These conditions are obviously necessary. Conversely, suppose that the conditions \(a\) , \(b\) , \(c\) are satisfied. On \(U(\mathrm{V})\) , the topology of pointwise convergence is identical to the topology of pointwise convergence in D (TVS, III, §3, No.~4, Prop. 5). Therefore the mapping \((s,x)\mapsto U(s)x\) of \(\mathbf{V}\times \mathbf{E}\) into E is continuous (GT, X, §2, No.~1, Cor. 3 of Prop. 1). Since \(U(s_0s)x = U(s_0)(U(s)x)\) for all \(s_0\in \mathbf{G}\) , \(s\in \mathbf{G}\) , \(x\in \mathbf{E}\) , one sees that \(U\) is continuous.

When \(\mathbf{G}\) is locally compact, the conditions \(a)\) and \(b)\) are equivalent to the condition:

\(a^{\prime})\) for every compact subset \(\mathbf{K}\) of \(\mathbf{G}\) , \(U(\mathbf{K})\) is equicontinuous.

\begin{enumerate}
\def\labelenumi{\arabic{enumi})}
\setcounter{enumi}{2}
\item
  Suppose that \(U\) is a continuous linear representation of \(G\) in \(E\) . For every \(s \in G\) , let \(\widehat{U}(s)\) be the continuous extension of \(U(s)\) to the completion \(\widehat{E}\) of \(E\) . Then \(\widehat{U}\) is a linear representation of \(G\) in \(\widehat{E}\) , satisfying conditions \(a)\) and \(c)\) of Remark 2, and also condition \(b)\) by GT, X, §2, No.~2, Prop. 4. Therefore \(\widehat{U}\) is a continuous linear representation of \(G\) in \(\widehat{E}\) .
\item
  When \(\mathbf{E}\) is a normed space, \(U\) is said to be isometric if \(\| U(s)\| = 1\) for every \(s\in \mathbf{G}\) . For this, it suffices that \(\| U(s)\| \leqslant 1\) for all \(s\in \mathbf{G}\) , because one then has
\end{enumerate}

\[
1 = \| 1 \| \leqslant \| U (s) \| \cdot \| U (s ^ {- 1}) \|,
\]

whence \(\| U(s)\| = \| U(s^{-1})\| = 1\) for all \(s\in G\) .

PROPOSITION 1. --- If G is a locally compact group and E is barreled, then every separately continuous linear representation U of G in E is continuous.

For every compact subset K of G, \(U(\mathrm{K})\) is compact for the topology of pointwise convergence (Remark 1), therefore is equicontinuous (TVS, III, §4, No.~2, Th. 1); one then applies Remark 2.

Lemma 1. --- Let G be a locally compact group, \(\rho\) a lower semi-continuous finite numerical function \(\geqslant 0\) on G such that \(\rho(st) \leqslant \rho(s)\rho(t)\) for all \(s, t \in \mathbf{G}\) . Then \(\rho\) is bounded above on every compact subset of G.

There exists a nonempty open subset U of G such that \(\rho\) is bounded above on U (GT, IX, §5, No.~4, Th. 2). Let K be a compact subset of G. Then K is covered by a finite number of sets \(s_{1}U,\ldots,s_{n}U\) . For every \(x\in U\) , one has \(\rho(s_{i}x)\leqslant\rho(s_{i})\rho(x)\) , therefore \(\rho\) is bounded above on the \(s_{i}U\) , hence on K.

Lemma 2. --- Let G be a topological group, U a linear representation of G in a normed space E, and A a dense subset of E. Assume that for every \(s \in G\) , \(U(s)\) is continuous, and that, for every \(x \in A\) , \(s \mapsto U(s)x\) is a continuous mapping of G into E. Then the function \(s \mapsto g(s) = \|U(s)\|\) on G is lower semi-continuous and satisfies \(g(st) \leqslant g(s)g(t)\) .

Let B be the unit ball of E. Then \(g(s) = \sup_{x \in \mathbf{B} \cap \mathbf{A}} \|U(s)x\|\) , and each function \(s \mapsto \|U(s)x\|\) is continuous on G, therefore g is lower semicontinuous. On the other hand,

\[
g (s t) = \| U (s) U (t) \| \leqslant \| U (s) \| \cdot \| U (t) \| = g (s) g (t).
\]

PROPOSITION 2. --- Let G be a locally compact group, U a linear representation of G in a normed space E. Let A be a dense subset of E. Assume that for every \(s \in G\) , \(U(s)\) is continuous and that, for every \(x \in A\) , \(s \mapsto U(s)x\) is a continuous mapping of G into E. Then U is continuous.

For, \(\|U(s)\|\) is bounded on every compact subset of G by Lemmas 1 and 2, and one then applies Remark 2.

